\section{Proofs of the Domain Transfers}
\label{app:domains}

\subsection{Regular Solutions and Countermodels}\label{sec:regular}
For a covariant system $\chi$ on a variable set $E$ of size $k_\chi$,
use graph states $(U,L)\in\mathcal P(E)^2$, with root state
$(D_\chi\mathord\uparrow z,D_\chi\mathord\downarrow z)$ for $z$.
Put
\begin{equation}\label{eq:g54}
 L_{z,\eta}=\{v:\Lo(\eta,v,z)\},\qquad
 U_{z,\eta}=\{v:\Up(\eta,z,v)\}.
\end{equation}
On letter $i$, replace $L$ by $\{v:\exists u\in L\ \Lo(i,v,u)\}$
and $U$ by $\{v:\exists u\in U\ \Up(i,u,v)\}$.
Let $F_-$ and $F_+$ be the roots of lower and upper constructor literals,
and $B_+$ the bottom-literal roots. Output
\[
 g(U,L)=\begin{cases}
 f,&L\cap F_-\ne\varnothing\land U\cap F_+\ne\varnothing,\\
 \bot,&L\cap F_-=\varnothing\land U\cap(B_+\cup F_+)\ne\varnothing,\\
 \top,&\text{otherwise}.
 \end{cases}
\]
Only constructor-labeled vertices have outgoing edges in the unfolding.

\begin{theorem}[Least-shape graph]\label{thm:shape-graph}
This finite graph unfolds at each root $z$ to $\sigma_{\mathrm{sh}}(z)$.
It is computable and has at most $2^{2k_\chi}$ states, with no consistency
hypothesis needed for the equality of trees.
\end{theorem}
\begin{proof}
The path rules (\cref{lem:path-rules}) at $\eta\mid i$ show
that the transition takes $(U_{z,\eta},L_{z,\eta})$ to
$(U_{z,\eta i},L_{z,\eta i})$. For the lower coordinate a lower
$\eta$-path from $u$ to $z$ composes with a lower $i$-path from $v$ to $u$;
for the upper coordinate paths compose in forward endpoint order.
The converse follows by factoring those judgments. The initial sets
retain full closure, so induction also covers the empty path.
Intersecting these sets with the three literal-root sets gives exactly
$\ell_f,u_f,u_\bot$ in \eqref{eq:g27}. Thus the output agrees with $h$.
A path exists in the unfolding iff all its proper-prefix outputs are $f$,
exactly the pruning predicate of the least shape. Hence labels and absent
paths both agree. Closure, transitions, and output tests are finite
computations on the $2^{2k_\chi}$ set pairs.
\end{proof}

\begin{theorem}[Regular countermodel transfer]\label{thm:regular-transfer}
Every unsafe word has a regular countermodel at the same word and side.
Consequently \eqref{eq:g7} holds, and so does \eqref{eq:sat-reg}.
\end{theorem}
\begin{proof}
Use the clash-free four-spine extension from \cref{thm:four-spine}.
Its least shape has a finite graph $G$ by \cref{thm:shape-graph}.
Selection also has a finite graph: use states $(r,s,b)$, initially
$(r_z,r_{\flip z},\operatorname{sg}(z))$, outputting the label at $r$
when $b=0$ and the dual of that at $s$ when $b=1$. On letter $i$, advance
both graph coordinates and update $b$ to $b\oplus c(i)$.
Common shape from \cref{lem:selector} ensures that retained coordinates
are both constructors or both extreme; stop at the extreme outputs.
Induction on retained paths gives \eqref{eq:g60}, including its
domain. At most $2|G|^2$ states suffice.
Decoding $H_0$ uses another polarity bit, initially zero, to flip output
labels. Restriction just selects roots, so these operations also preserve
graph finiteness.
The fixed solution and four prefix conditions proved in
\cref{thm:four-spine} therefore give a regular countermodel after decoding,
on either side. Conversely, a regular countermodel is an arbitrary one.
Any arbitrary countermodel supplies an unsafe word by \eqref{eq:g13},
proving \eqref{eq:g7}. The same construction without spines gives a regular
solution whenever a solution exists, which proves \eqref{eq:sat-reg}.
\end{proof}

\subsection{Finite Satisfiability and Least-Shape Depth}\label{sec:finite-sat}
We first work covariantly over $k_\chi$ variables, with the cycle clashes
\eqref{eq:g55} of \cref{sec:decision}: a nonempty word $\eta$ and variables
$a,b$ with $\Lo(\eta,a,a)$, $D_\chi(a,b)$, and $\Up(\eta,b,b)$.

\begin{theorem}[Covariant finite satisfiability]\label{thm:finite-sat}
Finite satisfiability is characterized by

\begin{equation}\label{eq:g56}
\Sat_{\mathrm{fin}}(\chi)
 \quad\Longleftrightarrow\quad
 \chi\text{ has neither a label clash nor a cycle clash}.
\end{equation}

\end{theorem}

\begin{proof}\emph{Necessity.} A label clash excludes every solution. For a cycle clash,
\eqref{eq:g23} gives
$\emb_{\bot}(\eta,\sigma(a))\le \sigma(a)\le \sigma(b)
\le\emb_{\top}(\eta,\sigma(b))$.
Iterating the two self-bounds gives the same sandwich along $\eta^h$
for every $h$. Along this path both bounding trees have constructor
labels at every proper prefix; the middle trees must therefore also
have constructors there. In particular, $\sigma(a)$ has a path of length
$h|\eta|$ for every $h$, contradicting finiteness.

\par\noindent\emph{Sufficiency.} By \cref{thm:sat}, absence of a label
clash makes $\sigma_{\mathrm{sh}}$ a solution. We prove that absence of a
cycle clash bounds its depth, independently of label consistency.
If a domain path of ${\sigma_{\mathrm{sh}}}(z)$ had length greater than $k_\chi^2+1$, its
prefix $\tau$ of length $k_\chi^2+1$ would have constructor label. By \eqref{eq:g27}
choose a lower and an upper path along $\tau$ from $z$, with endpoints
supporting their constructor labels. Factor both paths at every cut, giving
variables $a_i,b_i$, $0\le i\le |\tau|$, with

\begin{equation}\label{eq:g58}
\Lo(\tau[:i],a_i,z),\qquad
 \Up(\tau[:i],z,b_i),\qquad D_\chi(a_i,b_i).
\end{equation}

The last relation is cancellation \eqref{eq:g24}. To obtain simultaneous
factorizations with exact endpoints, induct on the positive path length:
factor its first letter, continue from the resulting variable, and absorb
the closure at each cut by transitivity. For length one retain the given
endpoints. This gives $a_0=b_0=z$ and the chosen terminal supports, with
$\Lo(\tau[i:j],a_j,a_i)$ and $\Up(\tau[i:j],b_i,b_j)$ for each
$i<j$, by composition. At $i=j$ use reflexivity. Thus no nontrivial
empty-path closure is being replaced by equality. Among the first $k_\chi^2+1$
pairs, two agree, at cuts $i<j$. The nonempty factor $\tau[i:j]$
then satisfies \eqref{eq:g55} at $(a_i,b_i)$, a contradiction. Therefore all
domains of ${\sigma_{\mathrm{sh}}}$ have depth at most $k_\chi^2+1$, and finite branching makes
them finite. This proves \eqref{eq:g56}.

For an effective cycle test, take vertices $(a,b)$ with $D_\chi(a,b)$
and an $i$-edge to $(a',b')$ when $\Lo(i,a',a)$ and $\Up(i,b,b')$.
Cancellation puts the target in the vertex set. A nonempty closed walk
composes to the two self-paths in \eqref{eq:g55}. Conversely, factor those
self-paths at every letter with their exact endpoints; at each cut,
transport their prefixes across $D_\chi(a,b)$ and cancel to obtain a
vertex pair. These pairs give a nonempty closed walk. Thus cycle clash is
exactly directed cyclicity in this finite graph. A simple cycle has at
most $k_\chi^2$ edges; equivalently it suffices to test words of lengths
$1$ through $k_\chi^2$. At arity zero, there are no edges.
\end{proof}

\begin{lemma}[Finite signed correspondence]\label{lem:finite-signed}
For variance, one still has the equivalence

\begin{equation}\label{eq:g59}
\Sat_{\mathrm{fin},c}(\phi)
 \quad\Longleftrightarrow\quad
 \Sat_{\mathrm{fin}}(\phi^{\pm}).
\end{equation}

\end{lemma}

\begin{proof}
Normalization preserves domains and proves the forward implication.
Conversely, a finite solution of $\psi$ excludes both clashes by
\cref{thm:finite-sat}. Its least shape has depth at most $(2k)^2+1$.
Selection by \cref{lem:selector} preserves that domain and gives a finite
fixed solution. Decoding by \cref{lem:signed} again preserves domains,
proving the reverse implication. Together with the two finite clash tests
this decides finite satisfiability for every variance.
\end{proof}

\subsection{Finite Countermodel Transfer}\label{sec:finite-transfer}
\begin{theorem}[Finite countermodel transfer]\label{thm:finite-transfer}
Suppose $\phi$ has a finite variance solution. Every arbitrary-tree
countermodel to $x\le y$ can then be replaced by a finite countermodel
at any chosen unsafe word and side of the original countermodel.
This proves \eqref{eq:g8}.
\end{theorem}

\begin{proof}
\par\noindent\emph{Spine extensions add no cycle clash.} This assertion holds for
the four-spine extension \eqref{eq:g28}, without a label-consistency
hypothesis. Assign old
variables rank zero, spine variables with index $j>0$ rank $j$, and
fillers rank $|w|+1$. Every new lower-constructor literal
has source children of strictly larger rank than its root; every new
upper-constructor literal has sink children of strictly larger rank.
The separation proof \eqref{eq:g30} still applies.

Induction on \eqref{eq:g22} shows that a lower path rooted at a source
ends at a source of at least its rank, strictly larger for a nonempty path.
At length zero separation forces identity. At the first letter, the closure
into the source forces the lower-literal root to be that source. Its literal
is new and its chosen child is a source of larger rank, so induction
applies to the remaining path. An upper path rooted at a sink satisfies
the dual assertion, using closure out of a sink. Neither kind can have a nonempty self-cycle there. In a
cycle clash \eqref{eq:g55}, $a$ cannot be a source and $b$ cannot be a sink.
If $a$ were a sink, $D(a,b)$ would force $a=b$, impossible for
the upper self-cycle; if $b$ were a source, the same argument contradicts
the lower self-cycle. Thus both are old. A path between old endpoints
cannot use a new constructor literal: its child would enter a source
or sink and could not return to an old endpoint. Closure and both path
judgments in \eqref{eq:g55} therefore reflect to the old system. Conversely, every
old cycle remains. This proves cycle-clash equivalence for the extension.

\par\noindent\emph{The selected witness is finite.}
Choose an unsafe word $w$ from the arbitrary countermodel.
By \cref{thm:four-spine}, its existing four-spine extension has no label
clash. Normalizing the assumed finite solution gives a finite solution of
$\psi$, so $\psi$ has no cycle clash. The rank argument
above excludes cycle clashes in the extension. Its least shape therefore
has depth at most $(2k+4|w|+4)^2+1$ by the depth proof of
\cref{thm:finite-sat}. The selector keeps precisely this finite domain.
It satisfies the whole extension and is fixed, so restriction supplies
all four conditions \eqref{eq:g18}; decoding gives a finite countermodel
at the same word, on either side.

The assumed finite solution is used only to exclude cycles in the base
constraints. The spine extension preserves this exclusion, so the
countermodel is reconstructed independently as a finite least-shape solution.
If no finite solution exists, finite entailment is vacuous. Otherwise, the construction
just proved takes every failure of arbitrary entailment to a finite
countermodel; conversely arbitrary entailment holds in every finite
assignment. These observations prove \eqref{eq:g8}.
\end{proof}

\subsection{Countermodel Bounds}\label{sec:countermodel-bounds}
For a chosen unsafe word $w$, put $k_\chi=2k+4|w|+4$. The regular
construction uses at most $2^{2k_\chi}$ least-shape graph states,
$2^{4k_\chi+1}$ selected graph states, and $2^{4k_\chi+2}$ states after
decoding. Under finite satisfiability, the finite construction has tree
depth at most $k_\chi^2+1$, not a bound of that size on the total number
of nodes. The finite search supplies $|w|\le H$ with $H$ from
\eqref{eq:bounds}, giving effective bounds in the input parameters.
